\chapter{Chapter06 — Small Bias Algorithm}
%%% generated by the blueprint pipeline from FABL.Chapter06.Constructions.SmallBiasAlgorithm
\section{Overview}

Declarations of the Lean module \texttt{FABL.Chapter06.Constructions.SmallBiasAlgorithm}.

\section{Declarations}

\begin{theorem}[Injectivity of the coefficient polynomial map]
\lean{FABL.binaryLowPolynomial_injective}
\label{FABL.binaryLowPolynomial_injective}
\leanok
\uses{FABL.binaryLowPolynomial, FABL.coeff_binaryLowPolynomial, FABL.𝔽₂}
\difficulty{4}
For every natural number $\ell$, the map sending a vector $a\in\mathbb F_2^\ell$ to
$\sum_{i=0}^{\ell-1}a_iX^i$ is injective.
\end{theorem}

\begin{theorem}[Injectivity of the encoding]
\lean{FABL.binaryAdjoinRootEncode_injective}
\label{FABL.binaryAdjoinRootEncode_injective}
\leanok
\uses{FABL.F₂Cube, FABL.binaryAdjoinRootEncode, FABL.binaryLowPolynomial_degree_lt, FABL.binaryLowPolynomial_injective, FABL.binaryMonicPolynomial_degree}
\difficulty{4}
Let $\ell$ be a nonnegative integer and let $m$ be a length-$\ell$ vector over $\mathbb
F_2$. Encoding length-$\ell$ vectors as their low-coefficient polynomials in the
quotient obtained by adjoining a root of $X^\ell+\sum_{i<\ell}m_iX^i$ is injective.
\end{theorem}

\begin{definition}[Binary adjoined-root linear map]
\lean{FABL.binaryAdjoinRootLinearMap}
\label{FABL.binaryAdjoinRootLinearMap}
\leanok
\uses{FABL.F₂Cube, FABL.binaryAdd, FABL.binaryAdjoinRootEncode, FABL.binaryAdjoinRootEncode_binaryAdd, FABL.binaryAdjoinRootEncode_binaryScaleAdd, FABL.binaryMonicPolynomial, FABL.binaryScaleAdd, FABL.𝔽₂}
For a natural number $\ell$ and a binary vector $m$ of length $\ell$, the binary cube of
length $\ell$ maps linearly over the field with two elements to the root-adjoined
algebra of the polynomial $X^\ell+\sum_{i=0}^{\ell-1}m_iX^i$, by the binary adjoin-root
encoding.
\end{definition}

\begin{definition}[Modular powers of binary coefficient vectors]
\lean{FABL.binaryPowMod}
\label{FABL.binaryPowMod}
\leanok
\uses{FABL.CertifiedBinaryFieldImplementation, FABL.F₂Cube, FABL.binaryMulMod, FABL.binaryOneCoefficients}
Let $\ell>0$, let $m$ be the low-coefficient vector of a monic irreducible polynomial of
degree $\ell$ over $\mathbb F_2$, and let $r\in\mathbb F_2^\ell$. Define a sequence of
coefficient vectors by $P_0=(1,0,\ldots,0)$ and $P_{e+1}=P_e r$ modulo that polynomial,
where multiplication is performed by the specified modular multiplication routine. The
value at exponent $e$ is $P_e$.
\end{definition}

\begin{theorem}[Powers under polynomial encoding]
\lean{FABL.binaryAdjoinRootEncode_binaryPowMod}
\label{FABL.binaryAdjoinRootEncode_binaryPowMod}
\leanok
\uses{FABL.CertifiedBinaryFieldImplementation, FABL.F₂Cube, FABL.adjoinRoot_mk_binaryMulMod, FABL.binaryAdjoinRootEncode, FABL.binaryLowPolynomial_binaryOneCoefficients, FABL.binaryMulMod, FABL.binaryPowMod}
\difficulty{4}
Let $\ell$ be a positive natural number, and let $m$ be a vector of length $\ell$ over
$\mathbb{F}_2$ such that the monic polynomial $X^\ell+\sum_{i=0}^{\ell-1}m_iX^i$ is
irreducible over $\mathbb{F}_2$. For every vector $r$ of length $\ell$ over
$\mathbb{F}_2$ and every natural number $e$, encoding the result of raising $r$ to the
$e$th power modulo this polynomial equals the $e$th power of the encoding of $r$ in the
quotient algebra.
\end{theorem}

\begin{definition}[Executable small bias generator]
\lean{FABL.executableSmallBiasGenerator}
\statementsource{Analysis-of-Boolean-Functions-by-Ryan-ODonnell}{
  pdf-ffbe9c7b3306-p153-b002,
  pdf-ffbe9c7b3306-p153-b003
}
\label{FABL.executableSmallBiasGenerator}
\leanok
\uses{FABL.CertifiedBinaryFieldImplementation, FABL.F₂Cube, FABL.binaryPowMod, FABL.f₂DotProduct}
For natural numbers $n$ and positive $\ell$, let a certified binary field implementation
of degree $\ell$ be given, and choose seed vectors $r,s\in(\mathbb{F}_2)^\ell$. The
generator returns a vector $x\in(\mathbb{F}_2)^n$ whose $i$th coordinate is the binary
dot product of $s$ with the coordinate vector representing $r^{i+1}$ modulo the
certified irreducible polynomial.
\end{definition}

\begin{definition}[List of generator outputs]
\lean{FABL.executableSmallBiasGeneratorList}
\statementsource{Analysis-of-Boolean-Functions-by-Ryan-ODonnell}{pdf-ffbe9c7b3306-p153-b003}
\label{FABL.executableSmallBiasGeneratorList}
\leanok
\uses{FABL.CertifiedBinaryFieldImplementation, FABL.F₂Cube, FABL.allBinaryVectors, FABL.executableSmallBiasGenerator}
For natural numbers $n$ and positive $\ell$, and given a certified binary field
implementation of degree $\ell$, form a list of vectors in $(\mathbb{F}_2)^n$ by
applying the executable small bias generator to every ordered pair of length-$\ell$
binary vectors, in the order determined by the Cartesian product of the lists of all
such vectors.
\end{definition}

\begin{definition}[Multiset of generator outputs]
\lean{FABL.executableSmallBiasGeneratorMultiset}
\statementsource{Analysis-of-Boolean-Functions-by-Ryan-ODonnell}{
  pdf-ffbe9c7b3306-p153-b002,
  pdf-ffbe9c7b3306-p153-b003
}
\label{FABL.executableSmallBiasGeneratorMultiset}
\leanok
\uses{FABL.CertifiedBinaryFieldImplementation, FABL.F₂Cube, FABL.executableSmallBiasGeneratorList}
For natural numbers $n$ and positive $\ell$, and given a certified binary field
implementation of degree $\ell$, take the multiset of vectors obtained from the list of
generator outputs for all ordered pairs of length-$\ell$ binary vectors.
\end{definition}

\begin{theorem}[Length of the small-bias generator list]
\lean{FABL.length_executableSmallBiasGeneratorList}
\label{FABL.length_executableSmallBiasGeneratorList}
\leanok
\uses{FABL.CertifiedBinaryFieldImplementation, FABL.executableSmallBiasGeneratorList, FABL.length_binaryVectorPairs}
\difficulty{1}
Let $n$ be a natural number and let $\ell$ be a positive natural number. For any
certified binary field implementation of degree $\ell$ (that is, one whose associated
monic polynomial is irreducible), the executable small-bias generator list has length
$2^{2\ell}$.
\end{theorem}

\begin{theorem}[Generator multiset cardinality]
\lean{FABL.executableSmallBiasGeneratorMultiset_card}
\proofsource{Analysis-of-Boolean-Functions-by-Ryan-ODonnell}{pdf-ffbe9c7b3306-p153-b003}
\label{FABL.executableSmallBiasGeneratorMultiset_card}
\leanok
\uses{FABL.CertifiedBinaryFieldImplementation, FABL.executableSmallBiasGeneratorMultiset, FABL.length_executableSmallBiasGeneratorList}
\difficulty{3}
Let $n$ be a natural number and let $\ell>0$. For any irreducible monic binary modulus
of degree $\ell$, the multiset of generator outputs, counting multiplicity over all
pairs of length-$\ell$ binary seeds, has cardinality $(2^\ell)^2$.
\end{theorem}

\begin{definition}[Generator output probability density]
\lean{FABL.executableSmallBiasGeneratorDensity}
\statementsource{Analysis-of-Boolean-Functions-by-Ryan-ODonnell}{
  pdf-ffbe9c7b3306-p153-b002,
  pdf-ffbe9c7b3306-p153-b003
}
\label{FABL.executableSmallBiasGeneratorDensity}
\leanok
\uses{FABL.CertifiedBinaryFieldImplementation, FABL.F₂Cube, FABL.ProbabilityDensity, FABL.ProbabilityDensity.uniformPushforward, FABL.executableSmallBiasGenerator}
For natural numbers $n$ and positive $\ell$, and given a certified binary field
implementation of degree $\ell$, define a probability density on $(\mathbb{F}_2)^n$ by
pushing forward the uniform distribution on pairs of length-$\ell$ binary vectors
through the executable small bias generator.
\end{definition}

\begin{definition}[Polynomial encoding a binary vector]
\lean{FABL.executableSmallBiasPolynomial}
\statementsource{Analysis-of-Boolean-Functions-by-Ryan-ODonnell}{pdf-ffbe9c7b3306-p153-b004}
\label{FABL.executableSmallBiasPolynomial}
\leanok
\uses{FABL.CertifiedBinaryFieldImplementation, FABL.F₂Cube, FABL.binaryMonicPolynomial, FABL.𝔽₂}
For a certified binary field implementation of degree $\ell$ and a vector
$\gamma\in(\mathbb{F}_2)^n$, define a polynomial over the field obtained by adjoining a
root of the implementation's irreducible modulus: its coefficient of $X^{i+1}$ is the
image of $\gamma_i$ under the natural field map, for each $i=0,\ldots,n-1$.
\end{definition}

\begin{theorem}[Evaluation of the generator polynomial]
\lean{FABL.executableSmallBiasPolynomial_eval}
\label{FABL.executableSmallBiasPolynomial_eval}
\leanok
\uses{FABL.CertifiedBinaryFieldImplementation, FABL.F₂Cube, FABL.binaryMonicPolynomial, FABL.executableSmallBiasPolynomial}
\difficulty{1}
Let $\ell,n\in\mathbb N$, let a certified binary field implementation of degree $\ell$
be given, and let $\gamma$ be a binary vector of length $n$. For every element $r$ of
the field extension defined by its irreducible modulus, evaluating the generator
polynomial associated with $\gamma$ at $r$ gives $\sum_{i=0}^{n-1}\gamma_i r^{i+1}$,
where each coefficient acts on the power by scalar multiplication.
\end{theorem}

\begin{theorem}[Executable small bias polynomial ne zero]
\lean{FABL.executableSmallBiasPolynomial_ne_zero}
\statementsource{Analysis-of-Boolean-Functions-by-Ryan-ODonnell}{pdf-ffbe9c7b3306-p153-b004}
\label{FABL.executableSmallBiasPolynomial_ne_zero}
\leanok
\uses{FABL.CertifiedBinaryFieldImplementation, FABL.F₂Cube, FABL.binaryMonicPolynomial, FABL.executableSmallBiasPolynomial, FABL.𝔽₂}
\difficulty{5}
% dataset-pipeline: stub -- prose pending build_informal_statements.py
Statement pending regeneration.
\end{theorem}

\begin{theorem}[Degree of the encoded polynomial]
\lean{FABL.executableSmallBiasPolynomial_natDegree_le}
\label{FABL.executableSmallBiasPolynomial_natDegree_le}
\leanok
\uses{FABL.CertifiedBinaryFieldImplementation, FABL.F₂Cube, FABL.executableSmallBiasPolynomial}
\difficulty{3}
Let $\ell,n\in\mathbb{N}$, let a certified irreducible binary field modulus of degree
$\ell$ be fixed, and let $\gamma$ be a length-$n$ vector over $\mathbb{F}_2$. Form the
polynomial over the quotient field by adjoining a root of that modulus, whose
coefficient of $X^{i+1}$ is the image of $\gamma_i$. Its natural degree is at most $n$.
\end{theorem}

\begin{definition}[Binary field power sum]
\lean{FABL.executableSmallBiasPowerSum}
\statementsource{Analysis-of-Boolean-Functions-by-Ryan-ODonnell}{
  pdf-ffbe9c7b3306-p153-b003,
  pdf-ffbe9c7b3306-p153-b004
}
\label{FABL.executableSmallBiasPowerSum}
\leanok
\uses{FABL.CertifiedBinaryFieldImplementation, FABL.F₂Cube, FABL.binaryPowMod}
For natural numbers $n$ and positive $\ell$, given a certified binary field
implementation of degree $\ell$, a vector $\gamma\in(\mathbb{F}_2)^n$, and a seed vector
$r\in(\mathbb{F}_2)^\ell$, define the length-$\ell$ binary vector obtained by summing,
over $i=0,\ldots,n-1$, the scalar multiple of the coordinate vector representing
$r^{i+1}$ modulo the certified irreducible polynomial, with scalar $\gamma_i$.
\end{definition}

\begin{theorem}[Evaluation of the power-sum polynomial]
\lean{FABL.binaryAdjoinRootEncode_executableSmallBiasPowerSum}
\label{FABL.binaryAdjoinRootEncode_executableSmallBiasPowerSum}
\leanok
\uses{FABL.CertifiedBinaryFieldImplementation, FABL.F₂Cube, FABL.binaryAdjoinRootEncode, FABL.binaryAdjoinRootEncode_binaryPowMod, FABL.binaryAdjoinRootLinearMap, FABL.binaryPowMod, FABL.executableSmallBiasPolynomial, FABL.executableSmallBiasPolynomial_eval, FABL.executableSmallBiasPowerSum}
\difficulty{3}
Let $n,\ell$ be nonnegative integers with $\ell>0$, let $m$ be a length-$\ell$ vector
over $\mathbb F_2$ for which $X^\ell+\sum_{j<\ell}m_jX^j$ is irreducible, and let
$\gamma\in\mathbb F_2^n$ and $r\in\mathbb F_2^\ell$. In the quotient obtained by
adjoining a root of that polynomial, encode the vector $\sum_{i<n}\gamma_i r^{i+1}$,
where powers and scalar multiples are computed in the length-$\ell$ modular
representation, by its low-coefficient polynomial. This encoding equals the value at the
encoding of $r$ of the polynomial $\sum_{i<n}\gamma_i X^{i+1}$, with coefficients mapped
into the quotient.
\end{theorem}

\begin{theorem}[Pairing with the small bias generator]
\lean{FABL.f₂DotProduct_executableSmallBiasGenerator}
\leanok
\uses{FABL.CertifiedBinaryFieldImplementation, FABL.F₂Cube, FABL.binaryPowMod, FABL.executableSmallBiasGenerator, FABL.executableSmallBiasPowerSum, FABL.f₂DotProduct}
\difficulty{4}
Let $n$ be a nonnegative integer, let $\ell>0$, and choose an irreducible monic binary
polynomial of degree $\ell$ to define the certified binary field implementation. For
binary vectors $\gamma$ of length $n$ and $r,s$ of length $\ell$, the dot product of
$\gamma$ with the generated length-$n$ vector whose $i$th coordinate is the pairing of
$s$ with $r^{i+1}$ equals the dot product of $s$ with the sum $\sum_{i=0}^{n-1}\gamma_i
r^{i+1}$, where the powers and sum use that field implementation.
\end{theorem}

\begin{theorem}[Walsh character identity for a finite field generator]
\lean{FABL.vectorWalshCharacter_executableSmallBiasGenerator}
\proofsource{Analysis-of-Boolean-Functions-by-Ryan-ODonnell}{pdf-ffbe9c7b3306-p153-b004}
\label{FABL.vectorWalshCharacter_executableSmallBiasGenerator}
\leanok
\uses{FABL.CertifiedBinaryFieldImplementation, FABL.F₂Cube, FABL.executableSmallBiasGenerator, FABL.executableSmallBiasPowerSum, FABL.f₂DotProduct_executableSmallBiasGenerator, FABL.vectorWalshCharacter, FABL.vectorWalshCharacter_apply}
\difficulty{2}
Let $n$ be a nonnegative integer and let $\ell>0$. Choose a monic irreducible polynomial
of degree $\ell$ over $\mathbb F_2$ to represent the binary extension field, and use its
coefficient representation to perform arithmetic in that field. For $\gamma\in\mathbb
F_2^n$ and $r,s$ in the represented field, form the vector whose $i$th coordinate is the
binary dot product of the coefficient vectors of $r^{i+1}$ and $s$. The Walsh character
indexed by this vector, evaluated at $x$ as prescribed by the generator, equals the
Walsh character indexed by the coefficient vector of $\sum_{i=0}^{n-1}\gamma_i r^{i+1}$,
evaluated at the coefficient vector of $s$.
\end{theorem}

\begin{theorem}[Walsh average of the polynomial generator]
\lean{FABL.expect_executableSmallBiasGenerator_fixed_first}
\proofsource{Analysis-of-Boolean-Functions-by-Ryan-ODonnell}{pdf-ffbe9c7b3306-p153-b004}
\label{FABL.expect_executableSmallBiasGenerator_fixed_first}
\leanok
\uses{FABL.CertifiedBinaryFieldImplementation, FABL.F₂Cube, FABL.executableSmallBiasGenerator, FABL.executableSmallBiasPowerSum, FABL.expect_vectorWalshCharacter, FABL.vectorWalshCharacter, FABL.vectorWalshCharacter_executableSmallBiasGenerator}
\difficulty{3}
Let $\ell>0$, and let a monic polynomial over $\mathbb F_2$ of degree $\ell$ be
irreducible. For $\gamma\in\mathbb F_2^n$ and $r\in\mathbb F_2^\ell$, average the Walsh
character indexed by $\gamma$ over the generator outputs obtained by choosing $s$
uniformly from $\mathbb F_2^\ell$. This average is $1$ if the sum of the vectors
$\gamma_i r^{i+1}$ in $\mathbb F_2^\ell$ is zero, and $0$ otherwise.
\end{theorem}

\begin{theorem}[Zeros of the executable power sum]
\lean{FABL.ncard_executableSmallBiasPowerSum_zero_le}
\proofsource{Analysis-of-Boolean-Functions-by-Ryan-ODonnell}{pdf-ffbe9c7b3306-p153-b004}
\label{FABL.ncard_executableSmallBiasPowerSum_zero_le}
\leanok
\uses{FABL.CertifiedBinaryFieldImplementation, FABL.F₂Cube, FABL.binaryAdjoinRootEncode, FABL.binaryAdjoinRootEncode_executableSmallBiasPowerSum, FABL.binaryAdjoinRootEncode_injective, FABL.binaryAdjoinRootEncode_zero, FABL.binaryMonicPolynomial, FABL.executableSmallBiasPolynomial, FABL.executableSmallBiasPolynomial_natDegree_le, FABL.executableSmallBiasPolynomial_ne_zero, FABL.executableSmallBiasPowerSum, FABL.ncard_rootSet_le_natDegree}
\difficulty{4}
Let $n$ be a natural number and let $\ell$ be a positive natural number. Choose a monic
binary polynomial of degree $\ell$ that is irreducible, and let $\gamma$ be a nonzero binary
vector of length $n$. Using the chosen polynomial to perform multiplication modulo it,
form the sum of $\gamma_i r^{i+1}$ over $0\le i<n$ for each binary vector $r$ of length
$\ell$. At most $n$ such vectors make this sum zero.
\end{theorem}

\begin{theorem}[Character expectation as a root probability]
\lean{FABL.executableSmallBiasGenerator_characterExpectation_eq_rootProbability}
\proofsource{Analysis-of-Boolean-Functions-by-Ryan-ODonnell}{pdf-ffbe9c7b3306-p153-b004}
\label{FABL.executableSmallBiasGenerator_characterExpectation_eq_rootProbability}
\leanok
\uses{FABL.CertifiedBinaryFieldImplementation, FABL.F₂Cube, FABL.executableSmallBiasGenerator, FABL.executableSmallBiasPowerSum, FABL.expect_executableSmallBiasGenerator_fixed_first, FABL.vectorWalshCharacter, FABL.𝔽₂}
\difficulty{4}
Let $n,\ell$ be natural numbers with $\ell>0$, choose an irreducible monic binary
modulus of degree $\ell$, and let $\gamma$ be any binary vector of length $n$. The
expected Walsh character of the generated output, when both length-$\ell$ seeds are
chosen uniformly, equals $2^{-\ell}$ times the number of length-$\ell$ binary vectors
$r$ for which $\sum_{i=0}^{n-1}\gamma_i r^{i+1}=0$ in the binary field defined by the
modulus.
\end{theorem}

\begin{theorem}[Nonnegativity of character expectation]
\lean{FABL.executableSmallBiasGenerator_characterExpectation_nonneg}
\proofsource{Analysis-of-Boolean-Functions-by-Ryan-ODonnell}{pdf-ffbe9c7b3306-p153-b004}
\label{FABL.executableSmallBiasGenerator_characterExpectation_nonneg}
\leanok
\uses{FABL.CertifiedBinaryFieldImplementation, FABL.F₂Cube, FABL.executableSmallBiasGenerator, FABL.executableSmallBiasGenerator_characterExpectation_eq_rootProbability, FABL.vectorWalshCharacter}
\difficulty{1}
Let $n,\ell$ be natural numbers with $\ell>0$, choose an irreducible monic binary
modulus of degree $\ell$, and let $\gamma$ be any binary vector of length $n$. The
expected Walsh character of the generated output, with both length-$\ell$ seeds chosen
uniformly, is nonnegative.
\end{theorem}

\begin{theorem}[Upper bound for character expectation]
\lean{FABL.executableSmallBiasGenerator_characterExpectation_le}
\proofsource{Analysis-of-Boolean-Functions-by-Ryan-ODonnell}{pdf-ffbe9c7b3306-p153-b004}
\label{FABL.executableSmallBiasGenerator_characterExpectation_le}
\leanok
\uses{FABL.CertifiedBinaryFieldImplementation, FABL.F₂Cube, FABL.executableSmallBiasGenerator, FABL.executableSmallBiasGenerator_characterExpectation_eq_rootProbability, FABL.ncard_executableSmallBiasPowerSum_zero_le, FABL.vectorWalshCharacter}
\difficulty{3}
Let $n,\ell$ be natural numbers with $\ell>0$, choose an irreducible monic binary
modulus of degree $\ell$, and let $\gamma$ be a nonzero binary vector of length $n$. The
expected Walsh character of the generated output, with both length-$\ell$ seeds chosen
uniformly, is at most $n/2^\ell$.
\end{theorem}

\begin{theorem}[Fourier coefficient of the binary generator density]
\lean{FABL.vectorFourierCoeff_executableSmallBiasGeneratorDensity}
\proofsource{Analysis-of-Boolean-Functions-by-Ryan-ODonnell}{pdf-ffbe9c7b3306-p153-b004}
\label{FABL.vectorFourierCoeff_executableSmallBiasGeneratorDensity}
\leanok
\uses{FABL.CertifiedBinaryFieldImplementation, FABL.F₂Cube, FABL.ProbabilityDensity.expectation, FABL.ProbabilityDensity.expectation_uniformPushforward, FABL.executableSmallBiasGeneratorDensity, FABL.executableSmallBiasGenerator_characterExpectation_eq_rootProbability, FABL.executableSmallBiasPowerSum, FABL.vectorFourierCoeff, FABL.vectorFourierCoeff_eq_expect, FABL.vectorWalshCharacter}
\difficulty{3}
Let $\ell>0$, and let a monic polynomial of degree $\ell$ over $\mathbb{F}_2$ be
specified by its lower coefficients and assumed irreducible. For every $n$-coordinate
binary vector $\gamma$, its vector Fourier coefficient under the distribution obtained
by choosing the generator’s two seeds uniformly is the number of roots
$r\in\mathbb{F}_2^\ell$ of the finite-field sum formed from the coordinates of $\gamma$
and the successive powers of the first seed, divided by $2^\ell$.
\end{theorem}

\begin{theorem}[Bias of the generator distribution]
\lean{FABL.executableSmallBiasGeneratorDensity_isBiased}
\proofsource{Analysis-of-Boolean-Functions-by-Ryan-ODonnell}{
  pdf-ffbe9c7b3306-p153-b002,
  pdf-ffbe9c7b3306-p153-b004
}
\label{FABL.executableSmallBiasGeneratorDensity_isBiased}
\leanok
\uses{FABL.CertifiedBinaryFieldImplementation, FABL.ProbabilityDensity.IsBiased, FABL.executableSmallBiasGeneratorDensity, FABL.executableSmallBiasGenerator_characterExpectation_eq_rootProbability, FABL.executableSmallBiasGenerator_characterExpectation_le, FABL.vectorFourierCoeff_executableSmallBiasGeneratorDensity}
\difficulty{4}
Let $n,\ell$ be natural numbers with $\ell>0$, and choose an irreducible monic binary
modulus of degree $\ell$. If $\varepsilon\in\mathbb R$ satisfies
$n/2^\ell\leq\varepsilon$, then the distribution of outputs obtained by choosing the two
length-$\ell$ binary seeds uniformly is $\varepsilon$-biased: for every nonzero binary
vector $\gamma$ of length $n$, the absolute value of its Fourier coefficient is at most
$\varepsilon$.
\end{theorem}

\begin{theorem}[Bias and size of the generator]
\lean{FABL.executableSmallBiasGenerator_core}
\proofsource{Analysis-of-Boolean-Functions-by-Ryan-ODonnell}{
  pdf-ffbe9c7b3306-p153-b003,
  pdf-ffbe9c7b3306-p153-b004
}
\label{FABL.executableSmallBiasGenerator_core}
\leanok
\uses{FABL.CertifiedBinaryFieldImplementation, FABL.ProbabilityDensity.IsBiased, FABL.executableSmallBiasGeneratorDensity, FABL.executableSmallBiasGeneratorDensity_isBiased, FABL.executableSmallBiasGeneratorMultiset, FABL.executableSmallBiasGeneratorMultiset_card}
\difficulty{4}
Let $n$ and $\ell$ be positive integers, and let a certified binary field implementation
of degree $\ell$ be given. If $\varepsilon>0$ and $n\leq \varepsilon 2^\ell\leq
4n/\varepsilon$, then the probability density generated by choosing two degree-$\ell$
binary vectors uniformly is $\varepsilon$-biased, and the cardinality of the generator
multiset is at most $16(n/\varepsilon)^2$.
\end{theorem}

\begin{definition}[Work for a binary power construction]
\lean{FABL.executableSmallBiasPowerWork}
\label{FABL.executableSmallBiasPowerWork}
\leanok
\uses{FABL.binaryConstructorTraversalWork, FABL.binaryMulModWork}
For coefficient-vector length $\ell$ and exponent bound $e$, define the work to be
$\ell$ plus the traversal work for the list of indices $0,\ldots,e-1$, with step cost
equal to one more than the work of a modular multiplication at length $\ell$.
\end{definition}

\begin{theorem}[Power work formula]
\lean{FABL.executableSmallBiasPowerWork_eq}
\label{FABL.executableSmallBiasPowerWork_eq}
\leanok
\uses{FABL.binaryConstructorTraversalWork_eq, FABL.binaryMulModWork, FABL.executableSmallBiasPowerWork}
\difficulty{1}
For all natural numbers $\ell$ and $e$, the work assigned to computing a power with
exponent $e$ is $\ell+e\bigl(W(\ell)+1\bigr)$, where $W(\ell)$ is the binary modular
multiplication work at degree $\ell$.
\end{theorem}

\begin{definition}[Work for binary row construction]
\lean{FABL.executableSmallBiasRowWork}
\label{FABL.executableSmallBiasRowWork}
\leanok
\uses{FABL.binaryConstructorTraversalWork, FABL.executableSmallBiasPowerWork}
For coefficient-vector length $\ell$, define the work recursively on the output
dimension: it is zero at dimension zero, and at dimension $n+1$ it is the work at
dimension $n$ plus the power-construction work for exponent $n+1$, plus twice the
traversal work over the indices $0,\ldots,\ell-1$, plus one.
\end{definition}

\begin{theorem}[Row work upper bound]
\lean{FABL.executableSmallBiasRowWork_le}
\label{FABL.executableSmallBiasRowWork_le}
\leanok
\uses{FABL.binaryConstructorTraversalWork_eq, FABL.binaryMulModWork, FABL.executableSmallBiasPowerWork_eq, FABL.executableSmallBiasRowWork}
\difficulty{4}
For all natural numbers $n$ and $\ell$, the work assigned to $n$ small-bias rows at
degree $\ell$ is at most $n\bigl(\ell+n(W(\ell)+1)+2\ell+1\bigr)$, where $W(\ell)$ is
the binary modular multiplication work at degree $\ell$.
\end{theorem}

\begin{definition}[Work for the executable small-bias construction]
\lean{FABL.executableSmallBiasConstructionWork}
\label{FABL.executableSmallBiasConstructionWork}
\leanok
\uses{FABL.allBinaryVectors, FABL.binaryConstructorTraversalWork, FABL.binaryFieldPreprocessingWork, FABL.executableSmallBiasRowWork}
For natural numbers $n$ and $\ell$, the work of the executable small-bias construction
is the sum of the binary-field preprocessing work at length $\ell$ and the
binary-constructor traversal work on the Cartesian product of the sets of binary vectors
of length $\ell$, with step cost equal to one plus the small-bias row work for
dimensions $n$ and $\ell$.
\end{definition}

\begin{theorem}[Small-bias construction work]
\lean{FABL.executableSmallBiasConstructionWork_eq}
\label{FABL.executableSmallBiasConstructionWork_eq}
\leanok
\uses{FABL.binaryConstructorTraversalWork_eq, FABL.binaryFieldPreprocessingWork, FABL.executableSmallBiasConstructionWork, FABL.executableSmallBiasRowWork, FABL.length_binaryVectorPairs}
\difficulty{1}
For all natural numbers $n$ and $\ell$, the executable small-bias construction work
equals the field preprocessing work for dimension $\ell$ plus $2^{2\ell}$ times one more
than the row work for $n$ rows and dimension $\ell$.
\end{theorem}

\begin{definition}[Executable small bias construction data]
\lean{FABL.ExecutableSmallBiasConstruction}
\statementsource{Analysis-of-Boolean-Functions-by-Ryan-ODonnell}{
  pdf-ffbe9c7b3306-p153-b002,
  pdf-ffbe9c7b3306-p153-b003
}
\label{FABL.ExecutableSmallBiasConstruction}
\leanok
\uses{FABL.ExecutableBinaryFieldModel, FABL.F₂Cube}
For a natural number $n$, an executable small bias construction consists of a positive
natural number $\ell$ specifying the field degree, an executable binary field model of
that degree with its arithmetic tables, and a list of length-$n$ binary output vectors
in the fixed ordered-pair seed order.
\end{definition}

\begin{definition}[Executable small-bias construction]
\lean{FABL.buildExecutableSmallBiasConstruction}
\statementsource{Analysis-of-Boolean-Functions-by-Ryan-ODonnell}{
  pdf-ffbe9c7b3306-p153-b002,
  pdf-ffbe9c7b3306-p153-b003
}
\label{FABL.buildExecutableSmallBiasConstruction}
\leanok
\uses{FABL.ExecutableSmallBiasConstruction, FABL.buildExecutableBinaryFieldModel, FABL.executableSmallBiasGeneratorList}
For natural numbers $n$ and $\varepsilon$ with $\varepsilon > 0$, define an executable small-bias
construction on $n$ coordinates whose field degree is $\varepsilon$, whose field model is a
certified binary field model of that degree with complete ordered addition and
multiplication tables, and whose output list consists of the generator outputs in the
fixed ordered-pair seed order.
\end{definition}

\begin{theorem}[Small-bias construction resource bounds]
\lean{FABL.buildExecutableSmallBiasConstruction_resource_bounds}
\label{FABL.buildExecutableSmallBiasConstruction_resource_bounds}
\leanok
\uses{FABL.binaryFieldPreprocessingWork, FABL.binaryFieldPreprocessingWork_le, FABL.binaryFieldRepresentationBits, FABL.binaryFieldRepresentationBits_le_preprocessingWork, FABL.buildExecutableBinaryFieldModel, FABL.buildExecutableSmallBiasConstruction, FABL.executableSmallBiasConstructionWork, FABL.executableSmallBiasConstructionWork_eq, FABL.executableSmallBiasRowWork, FABL.length_executableSmallBiasGeneratorList}
\difficulty{1}
Let $n$ be a natural number and let $\ell>0$. The constructed small-bias
construction has $2^{2\ell}$ outputs. Its field representation size is at most
its field preprocessing work, which is at most $2^{8(\ell+1)}$, and its
construction work equals the field preprocessing work plus $2^{2\ell}$ times one
more than the row work.
\end{theorem}

\begin{definition}[Small-bias input parameters]
\lean{FABL.SmallBiasInput}
\label{FABL.SmallBiasInput}
\leanok
A small-bias input consists of natural numbers $n$, $a$, and $b$, where $n>0$, $a>0$,
and $2a\le b$. The requested bias is encoded by the numerator $a$ and denominator $b$.
\end{definition}

\begin{definition}[Encoded bias]
\lean{FABL.SmallBiasInput.epsilon}
\label{FABL.SmallBiasInput.epsilon}
\leanok
\uses{FABL.SmallBiasInput}
For a small-bias input with numerator $a$ and denominator $b$, define its real bias
parameter to be $a/b$.
\end{definition}

\begin{definition}[Small-bias scale]
\lean{FABL.SmallBiasInput.scale}
\label{FABL.SmallBiasInput.scale}
\leanok
\uses{FABL.SmallBiasInput}
For a small-bias input with output dimension $n$, numerator $a$, and denominator $b$,
define its scale to be the ceiling of $nb/a$.
\end{definition}

\begin{definition}[Small-bias field degree]
\lean{FABL.SmallBiasInput.fieldDegree}
\label{FABL.SmallBiasInput.fieldDegree}
\leanok
\uses{FABL.SmallBiasInput, FABL.SmallBiasInput.scale}
For a small-bias input of scale $s$, define its field degree to be the larger of $1$ and
the least natural number $k$ such that $2^k\ge s$.
\end{definition}

\begin{theorem}[Positivity of the field degree]
\lean{FABL.SmallBiasInput.fieldDegree_pos}
\label{FABL.SmallBiasInput.fieldDegree_pos}
\leanok
\uses{FABL.SmallBiasInput, FABL.SmallBiasInput.fieldDegree}
\difficulty{1}
Let $n$, $a$, and $b$ be positive integers with $2a\le b$, and define the scale
$S=\lceil nb/a\rceil$ and field degree $D=\max(1,\lceil\log_2 S\rceil)$. Then $D>0$.
\end{theorem}

\begin{definition}[Deterministic small-bias algorithm]
\lean{FABL.deterministicSmallBiasAlgorithm}
\statementsource{Analysis-of-Boolean-Functions-by-Ryan-ODonnell}{
  pdf-ffbe9c7b3306-p153-b002,
  pdf-ffbe9c7b3306-p153-b003
}
\label{FABL.deterministicSmallBiasAlgorithm}
\leanok
\uses{FABL.ExecutableSmallBiasConstruction, FABL.SmallBiasInput, FABL.SmallBiasInput.fieldDegree, FABL.SmallBiasInput.fieldDegree_pos, FABL.buildExecutableSmallBiasConstruction}
Given small-bias input parameters $n,a,b$ with $n>0$, $a>0$, and $2a\le b$, the
algorithm returns the executable small-bias construction in dimension $n$ using the
input’s selected positive field degree.
\end{definition}

\begin{definition}[Deterministic small-bias multiset]
\lean{FABL.deterministicSmallBiasMultiset}
\statementsource{Analysis-of-Boolean-Functions-by-Ryan-ODonnell}{
  pdf-ffbe9c7b3306-p153-b002,
  pdf-ffbe9c7b3306-p153-b003
}
\label{FABL.deterministicSmallBiasMultiset}
\leanok
\uses{FABL.F₂Cube, FABL.SmallBiasInput, FABL.deterministicSmallBiasAlgorithm}
Given small-bias input parameters $n,a,b$ with $n>0$, $a>0$, and $2a\le b$, define the
multiset of generator outputs by taking the algorithm’s output list in dimension $n$ and
forgetting its order while retaining multiplicities.
\end{definition}

\begin{definition}[Deterministic small-bias density]
\lean{FABL.deterministicSmallBiasDensity}
\statementsource{Analysis-of-Boolean-Functions-by-Ryan-ODonnell}{
  pdf-ffbe9c7b3306-p153-b002,
  pdf-ffbe9c7b3306-p153-b004
}
\label{FABL.deterministicSmallBiasDensity}
\leanok
\uses{FABL.ProbabilityDensity, FABL.SmallBiasInput, FABL.deterministicSmallBiasAlgorithm, FABL.executableSmallBiasGeneratorDensity}
For natural numbers $n,a,b$ with $n>0$, $a>0$, and $2a\le b$, the deterministic
small-bias construction specifies a probability density on the binary cube of dimension
$n$.
\end{definition}

\begin{definition}[Deterministic small bias construction work]
\lean{FABL.deterministicSmallBiasWork}
\label{FABL.deterministicSmallBiasWork}
\leanok
\uses{FABL.SmallBiasInput, FABL.SmallBiasInput.fieldDegree, FABL.executableSmallBiasConstructionWork}
For a small-bias input, define its deterministic construction work to be the executable
construction's work count at the input's output dimension and field degree.
\end{definition}

\begin{theorem}[Cardinality of the deterministic output multiset]
\lean{FABL.deterministicSmallBiasMultiset_card}
\proofsource{Analysis-of-Boolean-Functions-by-Ryan-ODonnell}{pdf-ffbe9c7b3306-p153-b003}
\label{FABL.deterministicSmallBiasMultiset_card}
\leanok
\uses{FABL.SmallBiasInput, FABL.SmallBiasInput.fieldDegree, FABL.buildExecutableSmallBiasConstruction, FABL.deterministicSmallBiasAlgorithm, FABL.deterministicSmallBiasMultiset, FABL.length_executableSmallBiasGeneratorList}
\difficulty{3}
For accuracy data with $n>0$, positive numerator, and twice the numerator at most the
denominator, the deterministic small-bias output multiset has cardinality $(2^d)^2$,
where $d$ is the field degree determined by those data.
\end{theorem}

\begin{theorem}[Positivity of the bias]
\lean{FABL.SmallBiasInput.epsilon_pos}
\label{FABL.SmallBiasInput.epsilon_pos}
\leanok
\uses{FABL.SmallBiasInput, FABL.SmallBiasInput.epsilon}
\difficulty{4}
For a small-bias input with positive bias numerator, its bias parameter is positive.
\end{theorem}

\begin{theorem}[Bias is at most one half]
\lean{FABL.SmallBiasInput.epsilon_le_half}
\label{FABL.SmallBiasInput.epsilon_le_half}
\leanok
\uses{FABL.SmallBiasInput, FABL.SmallBiasInput.epsilon}
\difficulty{4}
For a small-bias input whose bias numerator is positive and twice its bias numerator is
at most its bias denominator, the bias parameter is at most $1/2$.
\end{theorem}

\begin{theorem}[Scale exceeds one]
\lean{FABL.SmallBiasInput.one_lt_scale}
\label{FABL.SmallBiasInput.one_lt_scale}
\leanok
\uses{FABL.SmallBiasInput, FABL.SmallBiasInput.scale}
\difficulty{4}
Let $n,a,b$ be positive integers with $2a\le b$, and define the scale by $s=\lceil
nb/a\rceil$. Then $1<s$.
\end{theorem}

\begin{theorem}[Field degree and logarithmic scale]
\lean{FABL.SmallBiasInput.fieldDegree_eq_clog}
\label{FABL.SmallBiasInput.fieldDegree_eq_clog}
\leanok
\uses{FABL.SmallBiasInput, FABL.SmallBiasInput.fieldDegree, FABL.SmallBiasInput.one_lt_scale, FABL.SmallBiasInput.scale}
\difficulty{3}
Let $n$, $a$, and $b$ be positive integers with $2a \le b$, and define the scale
$S=\lceil nb/a\rceil$. The field degree $\max(1,\lceil\log_2 S\rceil)$ equals
$\lceil\log_2 S\rceil$.
\end{theorem}

\begin{theorem}[Scale bounded by field size]
\lean{FABL.SmallBiasInput.scale_le_fieldSize}
\label{FABL.SmallBiasInput.scale_le_fieldSize}
\leanok
\uses{FABL.SmallBiasInput, FABL.SmallBiasInput.fieldDegree, FABL.SmallBiasInput.fieldDegree_eq_clog, FABL.SmallBiasInput.scale}
\difficulty{3}
Let $n,a,b$ be positive integers with $2a\le b$, and define $s=\lceil nb/a\rceil$ and
$d=\max(1,\lceil\log_2 s\rceil)$, where the ceiling-log expression denotes the least
natural number $k$ with $2^k\ge s$. Then $s\le 2^d$.
\end{theorem}

\begin{theorem}[Field size bound by twice the scale]
\lean{FABL.SmallBiasInput.fieldSize_le_two_scale}
\label{FABL.SmallBiasInput.fieldSize_le_two_scale}
\leanok
\uses{FABL.SmallBiasInput, FABL.SmallBiasInput.fieldDegree, FABL.SmallBiasInput.fieldDegree_eq_clog, FABL.SmallBiasInput.one_lt_scale, FABL.SmallBiasInput.scale}
\difficulty{4}
Let $n$, $a$, and $b$ be positive integers with $2a\le b$. Define $S=\lceil nb/a\rceil$
and $D=\max(1,\lceil\log_2 S\rceil)$. Then $2^D\le 2S$.
\end{theorem}

\begin{theorem}[Dimension bound by field size]
\lean{FABL.SmallBiasInput.dimension_le_epsilon_mul_fieldSize}
\label{FABL.SmallBiasInput.dimension_le_epsilon_mul_fieldSize}
\leanok
\uses{FABL.SmallBiasInput, FABL.SmallBiasInput.epsilon, FABL.SmallBiasInput.fieldDegree, FABL.SmallBiasInput.scale, FABL.SmallBiasInput.scale_le_fieldSize}
\difficulty{4}
For a small-bias input with dimension $n$, bias $\varepsilon$, and field degree $r$, the real
inequality $n\leq \epsilon 2^r$ holds.
\end{theorem}

\begin{theorem}[Scale numerator bound]
\lean{FABL.SmallBiasInput.scale_mul_numerator_le}
\label{FABL.SmallBiasInput.scale_mul_numerator_le}
\leanok
\uses{FABL.SmallBiasInput, FABL.SmallBiasInput.scale}
\difficulty{3}
Let $n,a,b$ be positive integers with $2a\le b$, and define $s=\lceil nb/a\rceil$. Then
$sa\le nb+a$.
\end{theorem}

\begin{theorem}[Field size times bias numerator bound]
\lean{FABL.SmallBiasInput.fieldSize_mul_numerator_le}
\label{FABL.SmallBiasInput.fieldSize_mul_numerator_le}
\leanok
\uses{FABL.SmallBiasInput, FABL.SmallBiasInput.fieldDegree, FABL.SmallBiasInput.fieldSize_le_two_scale, FABL.SmallBiasInput.scale, FABL.SmallBiasInput.scale_mul_numerator_le}
\difficulty{4}
Let $n$, $a$, and $b$ be positive integers with $2a\le b$. Define $S=\lceil nb/a\rceil$
and $D=\max(1,\lceil\log_2 S\rceil)$. Then $2^D a\le 4nb$.
\end{theorem}

\begin{theorem}[Field size bound in terms of dimension and bias]
\lean{FABL.SmallBiasInput.fieldSize_le_four_dimension_div_epsilon}
\label{FABL.SmallBiasInput.fieldSize_le_four_dimension_div_epsilon}
\leanok
\uses{FABL.SmallBiasInput, FABL.SmallBiasInput.epsilon, FABL.SmallBiasInput.fieldDegree, FABL.SmallBiasInput.fieldSize_mul_numerator_le}
\difficulty{4}
Let $n$, $a$, and $b$ be positive integers with $2a\le b$. Set $\varepsilon=a/b$, let
$S=\lceil nb/a\rceil$, and let $D=\max(1,\lceil\log_2 S\rceil)$. Then $2^D\le
4n/\varepsilon$.
\end{theorem}

\begin{theorem}[Dimension bounded by the scale]
\lean{FABL.SmallBiasInput.n_le_scale}
\label{FABL.SmallBiasInput.n_le_scale}
\leanok
\uses{FABL.SmallBiasInput, FABL.SmallBiasInput.scale}
\difficulty{4}
Let $n$, $a$, and $b$ be positive integers with $2a\le b$, and define $S=\lceil
nb/a\rceil$. Then $n\le S$.
\end{theorem}

\begin{theorem}[Upper bound for the field degree]
\lean{FABL.SmallBiasInput.fieldDegree_succ_le}
\label{FABL.SmallBiasInput.fieldDegree_succ_le}
\leanok
\uses{FABL.SmallBiasInput, FABL.SmallBiasInput.fieldDegree, FABL.SmallBiasInput.fieldSize_le_two_scale, FABL.SmallBiasInput.scale, FABL.nat_succ_le_two_pow_succ}
\difficulty{4}
Let $n$, $a$, and $b$ be positive integers with $2a\le b$, and define the scale
$S=\lceil nb/a\rceil$ and field degree $D=\max(1,\lceil\log_2 S\rceil)$. Then $D+1\le
4(S+1)$.
\end{theorem}

\begin{definition}[Polynomial budget]
\lean{FABL.SmallBiasInput.polynomialBudget}
\label{FABL.SmallBiasInput.polynomialBudget}
\leanok
\uses{FABL.SmallBiasInput, FABL.SmallBiasInput.scale}
For a small-bias input, define its polynomial budget to be $2^{17}(s+1)^8$, where $s$ is
the input's scale.
\end{definition}

\begin{theorem}[Work bounded by the polynomial budget]
\lean{FABL.deterministicSmallBiasWork_le_polynomialBudget}
\label{FABL.deterministicSmallBiasWork_le_polynomialBudget}
\leanok
\uses{FABL.SmallBiasInput, FABL.SmallBiasInput.fieldDegree, FABL.SmallBiasInput.fieldDegree_succ_le, FABL.SmallBiasInput.fieldSize_le_two_scale, FABL.SmallBiasInput.n_le_scale, FABL.SmallBiasInput.polynomialBudget, FABL.SmallBiasInput.scale, FABL.binaryAddWork, FABL.binaryArithmeticWork_le, FABL.binaryFieldPreprocessingWork, FABL.binaryFieldPreprocessingWork_le, FABL.binaryMulModWork, FABL.deterministicSmallBiasWork, FABL.executableSmallBiasConstructionWork_eq, FABL.executableSmallBiasRowWork, FABL.executableSmallBiasRowWork_le}
\difficulty{6}
For accuracy data with $n>0$, positive numerator, and twice the numerator at most the
denominator, the deterministic small-bias construction work is at most
$2^{17}(\text{scale}+1)^8$.
\end{theorem}

\begin{theorem}[Polynomial work bound for the construction]
\lean{FABL.deterministicSmallBiasWork_isBigO}
\label{FABL.deterministicSmallBiasWork_isBigO}
\leanok
\uses{FABL.SmallBiasInput, FABL.SmallBiasInput.polynomialBudget, FABL.SmallBiasInput.scale, FABL.deterministicSmallBiasWork, FABL.deterministicSmallBiasWork_le_polynomialBudget}
\difficulty{3}
As the scale of the small-bias input tends to infinity, the deterministic construction
work is $O((\text{scale}+1)^8)$. The inputs have positive dimension, positive accuracy
numerator, and twice the numerator at most the denominator.
\end{theorem}

\begin{theorem}[Guarantees for the deterministic small-bias algorithm]
\lean{FABL.deterministicSmallBiasAlgorithm_spec}
\proofsource{Analysis-of-Boolean-Functions-by-Ryan-ODonnell}{
  pdf-ffbe9c7b3306-p153-b002,
  pdf-ffbe9c7b3306-p153-b003,
  pdf-ffbe9c7b3306-p153-b004
}
\label{FABL.deterministicSmallBiasAlgorithm_spec}
\leanok
\uses{FABL.ProbabilityDensity.IsBiased, FABL.SmallBiasInput, FABL.SmallBiasInput.dimension_le_epsilon_mul_fieldSize, FABL.SmallBiasInput.epsilon, FABL.SmallBiasInput.epsilon_le_half, FABL.SmallBiasInput.epsilon_pos, FABL.SmallBiasInput.fieldDegree_pos, FABL.SmallBiasInput.fieldSize_le_four_dimension_div_epsilon, FABL.SmallBiasInput.polynomialBudget, FABL.buildExecutableSmallBiasConstruction, FABL.deterministicSmallBiasAlgorithm, FABL.deterministicSmallBiasDensity, FABL.deterministicSmallBiasMultiset, FABL.deterministicSmallBiasWork, FABL.deterministicSmallBiasWork_le_polynomialBudget, FABL.executableSmallBiasGeneratorMultiset, FABL.executableSmallBiasGenerator_core, FABL.smallBiasGeneratorDensity, FABL.smallBiasGeneratorMultiset, FABL.smallBiasGenerator_core}
\difficulty{3}
Let $n>0$ and let the rational accuracy have positive numerator with twice the numerator
at most its denominator. The algorithm has accuracy $0<\varepsilon\le 1/2$; its output
density is $\varepsilon$-biased and its output multiset has cardinality at most
$16(n/\varepsilon)^2$. The corresponding field-based generator density is also
$\varepsilon$-biased, its multiset has cardinality at most $16(n/\varepsilon)^2$, and
the algorithm's work is at most its polynomial budget.
\end{theorem}

\begin{theorem}[A small-bias generator with bounded size]
\lean{FABL.exists_smallBiasGenerator_of_real}
\proofsource{Analysis-of-Boolean-Functions-by-Ryan-ODonnell}{pdf-ffbe9c7b3306-p153-b002}
\label{FABL.exists_smallBiasGenerator_of_real}
\leanok
\uses{FABL.ProbabilityDensity.IsBiased, FABL.smallBiasGeneratorDensity, FABL.smallBiasGeneratorMultiset, FABL.smallBiasGenerator_core}
\difficulty{6}
For every integer $n\ge1$ and every $0<\varepsilon\le1/2$, there is a positive integer
$\ell$ such that the distribution generated by choosing two elements uniformly from the
binary extension field of degree $\ell$ is $\varepsilon$-biased on the $n$-dimensional
binary cube, and the multiset of generated strings has cardinality at most
$16(n/\varepsilon)^2$.
\end{theorem}

\begin{theorem}[Field degree for dyadic bias]
\lean{FABL.SmallBiasInput.fieldDegree_eq_of_dyadic}
\label{FABL.SmallBiasInput.fieldDegree_eq_of_dyadic}
\leanok
\uses{FABL.SmallBiasInput, FABL.SmallBiasInput.fieldDegree, FABL.SmallBiasInput.fieldDegree_eq_clog, FABL.SmallBiasInput.scale}
\difficulty{3}
Let $n$, $a$, and $b$ be positive integers with $2a\le b$, and define $S=\lceil
nb/a\rceil$ and the field degree $D=\max(1,\lceil\log_2 S\rceil)$. If $n=2^{d-r}$,
$a=1$, $b=2^r$, and $r\le d$, then $D=d$.
\end{theorem}

\begin{theorem}[Power-of-two small-bias parameters]
\lean{FABL.smallBiasGenerator_core_powerOfTwo}
\proofsource{Analysis-of-Boolean-Functions-by-Ryan-ODonnell}{
  pdf-ffbe9c7b3306-p153-b003,
  pdf-ffbe9c7b3306-p153-b004
}
\label{FABL.smallBiasGenerator_core_powerOfTwo}
\leanok
\uses{FABL.ProbabilityDensity.IsBiased, FABL.smallBiasGeneratorDensity, FABL.smallBiasGeneratorMultiset, FABL.smallBiasGenerator_core}
\difficulty{4}
Let $d$ and $b$ be integers with $d>0$ and $0\le b\le d$, and set $n=2^{d-b}$ and
$\varepsilon=2^{-b}$. Choose $r,s$ independently and uniformly from the binary extension
field of degree $d$, and form a binary vector of length $n$ whose coordinate indexed by
$i$ is the dot product over $\mathbb F_2$ of the basis-coordinate vectors of $r^{i+1}$
and $s$. The output distribution, viewed as a density relative to the uniform
distribution on the binary cube, has every nonzero Fourier coefficient of absolute value
at most $\varepsilon$; moreover, the multiset of outputs as $(r,s)$ ranges over all
field pairs has cardinality at most $16(n/\varepsilon)^2$.
\end{theorem}
